% !TeX program = pdflatex
% !BIB program = biber
%-----------------------------------------------------------------
%  Intermediate Project template
%  Warsaw University of Technology, Faculty of Power and
%  Aeronautical Engineering (MEiL)
%  --
%  Based on WUT-Thesis, copyleft by Artur M. Brodzki & Piotr Woźniak
%  https://github.com/ArturB/WUT-Thesis
%  Modified by Arseni Lysak, 2026. Licensed under GPL-3.0.
%-----------------------------------------------------------------

\documentclass[
    bindingoffset=5mm,  % Binding offset
    footnoteindent=3mm, % Footnote indent
    hyphenation=true    % Hyphenation turn on/off
]{src/wut-thesis}

\graphicspath{{img/}}            % Folder with figures
\addbibresource{references.bib} % Bibliography file

%-----------------------------------------------------------------
% Extra packages. The class already loads amsmath, amssymb,
% graphicx, listings, multirow, url, hyperref and biblatex.
%-----------------------------------------------------------------
\usepackage[T1]{fontenc}
\usepackage{ellipsis}
\usepackage[scaled=0.85]{beramono}  % Monospace font for code
\usepackage{xcolor}
\usepackage{subcaption}             % Sub-figures
\usepackage{float}                  % [H] float placement
\usepackage{booktabs}               % \toprule, \midrule, \bottomrule
\usepackage{siunitx}                % Units: \SI{10}{\meter\per\second}
\usepackage{pdflscape}              % Landscape pages
\usepackage[section]{placeins}      % Keep floats inside their section
\usepackage[numbered,framed]{matlab-prettifier} % MATLAB listings
\lstset{
	inputencoding=utf8,
	extendedchars=true,
	literate={θ}{{$\theta$}}1       % Allow θ in MATLAB code
}
% Style for other code: \begin{lstlisting}[style=code, language=Python]
\definecolor{codegreen}{rgb}{0,0.6,0}
\definecolor{codegray}{rgb}{0.5,0.5,0.5}
\definecolor{codepurple}{rgb}{0.58,0,0.82}
\lstdefinestyle{code}{
	basicstyle=\ttfamily\footnotesize,
	keywordstyle=\color{blue},
	commentstyle=\color{codegreen},
	stringstyle=\color{codepurple},
	numbers=left,
	numberstyle=\tiny\color{codegray},
	breaklines=true,
	showstringspaces=false,
	tabsize=4
}
% Optional: nicer highlighting with minted, needs Python (see README)
% \usepackage{minted}
% \setminted{breaklines=true}

\raggedbottom
\emergencystretch=3em
\setcounter{tocdepth}{2}

%-----------------------------------------------------------------
% Title page
%-----------------------------------------------------------------
\facultymeil % \facultymeil or \facultyeiti
\langeng     % \langeng or \langpol

\institute{Aeronautics and Applied Mechanics} % Printed as "Institute of ..."
\worktype{Intermediate Project}              % Type of work, large and bold
\fieldofstudy{Aerospace Engineering}         % Free-text line under the type of work
\title{Title of the Intermediate Project}
\author{Name Surname 123456}                 % Name and student ID
\supervisor{Supervisor's name}               % Printed under "Project supervisor:"
\date{\the\year}

\begin{document}
	\maketitle
	\newpage

	\tableofcontents
	\newpage

	\section{Introduction}

	Motivation, the problem and the goal of the project.
	Cite sources with \verb|\cite|, e.g.~\cite{lamport1994latex, wutthesis}.

	\section{Related work}

	Overview of existing methods and their limitations.

	\section{Methodology}

	Equations are numbered and referenced with \verb|\eqref|, e.g.~Equation~\eqref{eq:newton}:
	\begin{equation} \label{eq:newton}
		\vec{F} = m \vec{a}
	\end{equation}
	Units are typeset with \texttt{siunitx}, e.g.~\SI{9.81}{\meter\per\second\squared}.

	Code is typeset with \texttt{listings}, as in Listing~\ref{lst:example}:
	\begin{lstlisting}[style=code, language=Python, caption={Example Python listing}, label={lst:example}]
def acceleration(force, mass):
    # Newton's second law
    return force / mass
	\end{lstlisting}

	\section{Results}

	Figure~\ref{fig:example} and Table~\ref{tab:example} show how to include figures and tables.
	Put your images in the \texttt{img} folder.

	\begin{figure}[htbp]
		\centering
		\includegraphics[width=0.6\textwidth]{example-image}
		\caption{Example figure.}
		\label{fig:example}
	\end{figure}

	\begin{table}[htbp]
		\centering
		\caption{Example table.}
		\label{tab:example}
		\begin{tabular}{l S[table-format=2.1] S[table-format=2.2]}
			\toprule
			\textbf{Case} & {\textbf{Accuracy (\%)}} & {\textbf{Score}} \\
			\midrule
			A & 54.2 & 57.72 \\
			B & 12.0 & 8.50 \\
			\bottomrule
		\end{tabular}
	\end{table}

	\section{Conclusions}

	Summary of the results and directions for future work.

	\printbibliography

\end{document}
